\chapter{I/O e streaming: lendo arquivos grandes sem sofrer}

\section{O que e I/O?}

I/O significa \textit{Input/Output}. Em termos simples, e tudo que envolve entrada ou saida de dados fora da CPU. Ler um arquivo, escrever em disco, consultar banco e esperar uma resposta de rede sao exemplos de I/O.

No seu projeto, as operacoes de I/O mais importantes sao:

\begin{itemize}
  \item ler \code{export.xml};
  \item ler arquivos \code{.gpx};
  \item inserir dados no Postgres;
  \item escrever logs;
  \item consultar tabelas de insights.
\end{itemize}

\section{Por que o XML gigante muda a estrategia?}

Se o arquivo tem 280 MB ou 628 MB, carregar tudo em memoria pode ser ruim. Alem do texto bruto, bibliotecas como BeautifulSoup criam objetos Python para representar a arvore XML inteira. Isso pode multiplicar o uso de memoria.

\begin{warningbox}{Evite no arquivo gigante}
\begin{lstlisting}[style=devbutterpython]
content = file.read_text(encoding="utf-8")
soup = BeautifulSoup(content, "xml")
\end{lstlisting}
\end{warningbox}

BeautifulSoup ainda pode ser util para arquivos pequenos, testes e inspecoes pontuais. Mas nao deve ser o parser principal do export gigante.

\section{Streaming com iterparse}

A alternativa e usar \code{xml.etree.ElementTree.iterparse}. Ele percorre o XML elemento por elemento.

\begin{lstlisting}[style=devbutterpython]
from pathlib import Path
from collections.abc import Iterator
import xml.etree.ElementTree as ET


def clean_tag(tag: str) -> str:
    if "}" in tag:
        return tag.split("}", 1)[1]
    return tag


def stream_record_elements(xml_file: Path) -> Iterator[ET.Element]:
    context = ET.iterparse(xml_file, events=("end",))

    for _, element in context:
        tag = clean_tag(element.tag)

        if tag == "Record":
            yield element

        element.clear()
\end{lstlisting}

\section{Explicando cada nomezinho}

\begin{itemize}
  \item \code{Path}: objeto da biblioteca \code{pathlib} para lidar com caminhos de arquivo de forma mais segura que string.
  \item \code{Iterator}: tipo que indica que a funcao retorna valores aos poucos.
  \item \code{ET}: apelido comum para \code{xml.etree.ElementTree}.
  \item \code{iterparse}: parser incremental de XML.
  \item \code{events=("end",)}: processa o elemento quando ele fecha, ou seja, quando ja temos seus atributos e filhos.
  \item \code{yield}: retorna um item sem encerrar a funcao; cria um generator.
  \item \code{element.clear()}: limpa o elemento da memoria depois de processar.
\end{itemize}

\section{Complexidade}

Se o XML tem \(n\) elementos, o parser passa por cada elemento uma vez:

\[
T(n) = O(n)
\]

A memoria tende a ser quase constante:

\[
M(n) \approx O(1)
\]

Isso nao significa literalmente zero crescimento, mas significa que voce nao precisa manter o XML inteiro em memoria.

\section{Primeiro exercicio}

Crie um XML pequeno em \code{tests/fixtures/small_export.xml} e teste se sua stream encontra apenas tags \code{Record}.

\begin{lstlisting}[style=devbutterpython]
def test_stream_record_elements_returns_records_only(tmp_path):
    xml_file = tmp_path / "small_export.xml"
    xml_file.write_text("""
    <HealthData>
        <Record type="HKQuantityTypeIdentifierHeartRate" value="90" />
        <Workout workoutActivityType="HKWorkoutActivityTypeRunning" />
    </HealthData>
    """)

    records = list(stream_record_elements(xml_file))

    assert len(records) == 1
    assert records[0].attrib["type"] == "HKQuantityTypeIdentifierHeartRate"
\end{lstlisting}
